\subsection{NOWHERE DENSE SETS}

DEFINITION 1. A subset A of a topological space X is said to be nowhere dense if its closure has no interior points.

Equivalently, A is nowhere dense in X if and only if the exterior of A is dense in X.

A closed set A is nowhere dense if and only if it has no interior points; that is, if and only if it coincides with its frontier. An arbitrary subset A is nowhere dense if and only if the closure of A is nowhere dense. Every subset of a nowhere dense set is nowhere dense.

Examples. 1) The empty subset of X is nowhere dense. In a Hausdorff space X, a set consisting of a single point is nowhere dense if and only if the point is not isolated in X. A dense subset is never nowhere dense (unless \(X = \varnothing\) ).

\begin{enumerate}
\def\labelenumi{\arabic{enumi})}
\setcounter{enumi}{1}
\item
  The frontier of an open or of a closed set is always nowhere dense.
\item
  In the space \(\mathbf{R}^n\) , every linear subspace of dimension \(p < n\) is a nowhere dense set (Chapter VI, § 1, no. 4, Proposition 2).
\end{enumerate}

Remark. The frontier of an arbitrary subset need not be nowhere dense: for example, if A and \(\mathbb{C}\) A are both dense sets, then the frontier of A is the whole space.

PROPOSITION I. The union of a finite number of nowhere dense sets is nowhere dense.

It is enough to show that the union of two nowhere dense sets A, B is nowhere dense, and without loss of generality we may assume that A and B are closed. The proposition is then equivalent to saying that the intersection of two dense open sets CA, CB is dense. Now if U is a non-empty open set, then \(U \cap CA\) is open and non-empty, hence

\[
\left(\mathrm{U} \cap \mathbb {C A}\right) \cap \mathbb {C B} = \mathrm{U} \cap \left(\mathbb {C A} \cap \mathbb {C B}\right)
\]

is open and non-empty.

Let Y be a subspace of a topological space X. A subspace A of Y is said to be nowhere dense relative to Y if A is nowhere dense when considered as a subset of the topological space Y.

PROPOSITION 2. Let Y be a subspace of a topological space X, and let A be a subset of Y. If A is nowhere dense relative to Y, then A is nowhere dense relative to X. Conversely, if Y is open in X and A is nowhere dense relative to X, then A is nowhere dense relative to Y.

Suppose that A is nowhere dense relative to Y. If the closure \(\overline{A}\) of A in X contains a non-empty open set U, then \(U \cap A\) is not empty (by the definition of closure); hence \(U \cap Y\) is a non-empty open set relative to Y, and is contained in the closure \(\overline{A} \cap Y\) of A with respect to Y, which is contrary to the hypothesis.

Now suppose that Y is open in X and that \(A \subset Y\) is nowhere dense relative to X. If U is open in Y and is not empty, then U is open in X and therefore contains a non-empty set V which is open in X (and a fortiori in Y) and does not meet A; hence A is nowhere dense relative to Y.

The second part of Proposition 2 is clearly not valid if Y is not open in X; consider, e.g., the situation where \(Y \neq \varnothing\) is nowhere dense in X, and A = Y.

